\section{Bucle de búsqueda común}
\label{sec:base}

La búsqueda adaptativa en vecindarios grandes (ALNS)~\cite{ropke2006,pisinger2007} mejora una
solución destruyéndola parcialmente y reparándola en cada iteración. El
Algoritmo~\ref{alg:alns} reproduce el método \cod{ALNS::solve}, común a los dos métodos de la
tesis. Las llamadas \fn{Seleccionar} y \fn{Aprender} son los dos únicos puntos de variación
(Secciones~\ref{sec:clasico} y~\ref{sec:ql}).

La búsqueda se organiza en \emph{dos fases}, siguiendo el esquema de Ropke y
Pisinger~\cite{ropke2006}: la fase~1 minimiza el número de rutas y la fase~2 minimiza la
distancia con la flota obtenida. Ambas ejecutan exactamente el mismo cuerpo de iteración; lo
que cambia es la solución de la que se parte y las condiciones de corte
(Sección~\ref{sec:fases}). En el pseudocódigo, $\bar{x}$ denota la última solución completa
conocida y $T_{\mathrm{fin}}$ la última iteración, que se fija al comenzar la fase~2.

\begin{algorithm}[H]
\caption{Bucle ALNS común en dos fases (\cod{ALNS::solve})}
\label{alg:alns}
\Entrada{solución inicial $x_0$, presupuesto de iteraciones $T_{\max}$}
\Salida{mejor solución encontrada $\xbest$}
$\bar{x}\gets x_0$;\quad $\textit{fase1}\gets\text{verdadero}$;\quad $T_{\mathrm{fin}}\gets 2\,T_{\max}$\;
$T_0\gets -\,w\,D(x_0)/\ln 0.5$;\quad $K_{\min}\gets\bigl\lceil\sum_{i}\mu_i/Q\bigr\rceil$\;
\fn{IniciarEtapa}($0$) \Com*[r]{Algoritmo~\ref{alg:etapa}: fija $x$, $\xbest$, $T$ y los contadores}
$t\gets 1$\;
\Mientras{$t\le T_{\mathrm{fin}}$}{
  $x'\gets x$ \Com*[r]{copia de trabajo}
  $(j,k)\gets$ \fn{Seleccionar}($t$) \Com*[r]{destrucción $j$, reparación $k$}
  \fn{AplicarOperadores}($x'$, $j$, $k$) \Com*[r]{Algoritmo~\ref{alg:aplicar}}
  $b_{\mathrm{ant}}\gets|\Uns(\xbest)|$\;
  $o.\gbest\gets G(\xbest,x')$;\quad $o.\gcur\gets G(x,x')$ \Com*[r]{ganancias relativas, ec.~\eqref{eq:ganancia}}
  $o.\textit{nuevoMejor}\gets f(x')<f(\xbest)$\;
  $o.\textit{mejoró}\gets o.\textit{nuevoMejor}\ \vee\ f(x')<f(x)$\;
  $o.\textit{aceptado}\gets o.\textit{mejoró}\ \vee\ $\fn{Aceptar}($f(x')$, $f(x)$, $T$)\;
  \lSi{$o.\textit{aceptado}$}{$x\gets x'$}
  \lSi{$o.\textit{nuevoMejor}$}{$\xbest\gets x'$}
  $\textit{estanc}\gets 0$ si $o.\textit{nuevoMejor}$; si no, $\textit{estanc}+1$\;
  \Si(\Com*[f]{reinicio por estancamiento}){$\textit{estanc}\ge L_{\mathrm{est}}$}{
    $x\gets\xbest$;\quad $T\gets\rho\,T_0$;\quad $\textit{estanc}\gets 0$\;
  }
  \fn{Aprender}($t$, $o$)\;
  $T\gets c\,T$ \Com*[r]{enfriamiento geométrico}
  \Si(\Com*[f]{control de la fase de rutas}){$\textit{fase1}$}{
    $b\gets|\Uns(\xbest)|$\;
    $\textit{sinBanco}\gets 0$ si $b<b_{\mathrm{ant}}$; si no, $\textit{sinBanco}+1$\;
    \uSi(\Com*[f]{ruta eliminada: se intenta con la siguiente}){$b=0$}{
      $\bar{x}\gets\xbest$;\quad \fn{IniciarEtapa}($t$)\;
    }
    \SiNoSi(\Com*[f]{se abandona el intento}){$t\ge T_{\max}$ \O $(b\ge B$ \Y $\textit{sinBanco}\ge L_{\mathrm{rut}})$}{
      $\textit{fase1}\gets\text{falso}$;\quad \fn{IniciarEtapa}($t$)\;
    }
  }
  $t\gets t+1$\;
}
\Devolver $\xbest$\;
\end{algorithm}

\begin{algorithm}[H]
\caption{\fn{IniciarEtapa} (función local \cod{startStage} de \cod{ALNS::solve})}
\label{alg:etapa}
\Entrada{iteración $t$ en la que comienza la etapa}
$\textit{fase1}\gets\textit{fase1}\ \wedge\ t<T_{\max}\ \wedge\ K(\bar{x})>K_{\min}$\;
\Si(\Com*[f]{comienza la fase 2}){$\neg\,\textit{fase1}$}{
  $T_{\mathrm{fin}}\gets t+T_{\max}$\;
  \fn{AlIniciarFase2}($t$) \Com*[r]{aviso a la subclase (\cod{onDistancePhase})}
}
$x\gets\bar{x}$\;
\Si(\Com*[f]{fase 1: una ruta menos}){$\textit{fase1}$}{
  eliminar de $\Rset(x)$ las rutas vacías\;
  $\pi\gets$ ruta de $x$ con menos clientes\;
  mover todos los clientes de $\pi$ a $\Uns(x)$;\quad eliminar $\pi$ de $\Rset(x)$\;
}
$\xbest\gets x$;\quad $T\gets T_0$;\quad $\textit{estanc}\gets 0$;\quad $\textit{sinBanco}\gets 0$\;
\end{algorithm}

\begin{table}[H]
\centering
\small
\caption{Constantes del bucle común (\cod{alns.cpp}, \cod{operators.cpp}, \cod{solution.cpp}).}
\label{tab:const-base}
\begin{tabularx}{\textwidth}{@{}l c c L@{}}
\toprule
Constante & Símbolo & Valor & Función \\
\midrule
\cod{UNASSIGNED\_COST} & -- & $10^6$ & Peso de un cliente sin asignar en $f$ \\
\cod{VEHICLE\_COST} & -- & $10^4$ & Peso de un vehículo en $f$ \\
\cod{TEMP\_SCALE} & $w$ & $0.2$ & Empeoramiento relativo aceptado con probabilidad $0.5$ en $T_0$ \\
\cod{COOLING\_RATE} & $c$ & $0.9995$ & Factor de enfriamiento por iteración \\
\cod{STAGNATION\_LIMIT} & $L_{\mathrm{est}}$ & $1500$ & Iteraciones sin nuevo mejor que disparan el reinicio \\
\cod{REHEAT\_FACTOR} & $\rho$ & $0.5$ & Fracción de $T_0$ a la que se recalienta \\
\cod{ROUTE\_PHASE\_BANK} & $B$ & $5$ & Clientes sin asignar a partir de los cuales puede abandonarse un intento de la fase~1 \\
\cod{ROUTE\_PHASE\_STALL} & $L_{\mathrm{rut}}$ & $2000$ & Iteraciones sin reducir los clientes sin asignar que disparan ese abandono \\
\cod{Q\_MIN\_RATIO} / \cod{Q\_MAX\_RATIO} & -- & $0.10$ / $0.40$ & Fracción mínima y máxima de clientes extraídos \\
\cod{Q\_MIN\_CAP} / \cod{Q\_MAX\_CAP} & -- & $30$ / $60$ & Topes absolutos de esas dos cotas \\
\cod{Q\_MIN\_FLOOR} & -- & $4$ & Cota inferior absoluta de $q$ \\
\cod{NOISE\_PROB} & -- & $0.5$ & Probabilidad de que la reparación use costos con ruido \\
\bottomrule
\end{tabularx}
\end{table}

\subsection{Las dos fases de la búsqueda}
\label{sec:fases}

\paragraph{Flota fija.} Los operadores de reparación \emph{nunca abren rutas nuevas}: insertan
los clientes pendientes en las rutas que ya existen en la solución y, si alguno no cabe en
ninguna, lo dejan en $\Uns$ (Sección~\ref{sec:reparacion}). El número de rutas de una
solución solo puede cambiar, por tanto, cuando el bucle común retira una de forma explícita.
Esta es la diferencia de fondo con un ALNS de una sola fase, en el que la reparación abre
rutas a demanda y la reducción de flota se confía a la función de costo: aquí la flota es un
dato de cada etapa y la búsqueda debe encajar a todos los clientes en ella.

\paragraph{Fase 1: minimización de rutas.} Se compone de una sucesión de \emph{etapas}. Cada
etapa parte de la última solución completa $\bar{x}$, le quita la ruta con menos clientes y
deja esos clientes sin asignar (Algoritmo~\ref{alg:etapa}). La lista $\Uns$ actúa como el
\emph{banco de clientes} de~\cite{ropke2006}: los clientes pendientes viajan con la solución
de una iteración a la siguiente y cada reparación intenta colocarlos junto con los $q$ recién
extraídos. Como $f$ penaliza con $10^6$ cada cliente sin asignar, el ALNS ordinario, sin
ninguna modificación, trabaja para vaciar el banco. La etapa termina de una de estas tres
formas:
\begin{itemize}
\item \textbf{Éxito.} La mejor solución de la etapa queda completa ($|\Uns(\xbest)|=0$): se ha
eliminado una ruta. Esa solución pasa a ser $\bar{x}$ y comienza otra etapa, con una ruta
menos.
\item \textbf{Abandono por estancamiento.} Quedan al menos $B=5$ clientes sin asignar en la
mejor solución y han pasado $L_{\mathrm{rut}}=2000$ iteraciones sin que ese número disminuya.
Es el criterio de corte de~\cite{ropke2006}: si el banco es grande y no se reduce, la ruta no
es eliminable con un esfuerzo razonable. Con menos de $5$ clientes pendientes el intento no se
abandona por esta vía y continúa hasta agotar el presupuesto.
\item \textbf{Presupuesto agotado.} Se alcanza la iteración $T_{\max}$.
\end{itemize}
En los dos últimos casos el intento en curso se descarta y la fase~1 termina. También termina,
sin iniciar una etapa nueva, cuando $\bar{x}$ alcanza la cota inferior trivial de vehículos
\begin{equation}
K_{\min}=\Bigl\lceil\frac{1}{Q}\sum_{i=1}^{n}\mu_i\Bigr\rceil ,
\label{eq:kmin}
\end{equation}
ya que con menos rutas no hay capacidad para toda la demanda. Si la solución inicial ya cumple
$K(x_0)=K_{\min}$, la fase~1 no se ejecuta.

\paragraph{Fase 2: minimización de distancia.} Parte de $\bar{x}$, la solución completa con
menos rutas que encontró la fase~1 (no del intento fallido), y ejecuta siempre $T_{\max}$
iteraciones: si la fase~1 terminó en la iteración $t_1\le T_{\max}$, la búsqueda acaba en
$T_{\mathrm{fin}}=t_1+T_{\max}$. La flota sigue fija, de modo que el número de vehículos no
puede superar el número de rutas de $\bar{x}$; sí puede bajar si una destrucción vacía una
ruta y la reparación no vuelve a usarla, lo que $f$ premia con $10^4$. Una candidata
incompleta cuesta al menos $10^6$ más que la solución actual, por lo que en esta fase es
rechazada en la práctica (Sección~\ref{sec:aceptacion}) y nunca es un nuevo mejor: la solución
devuelta es completa.

\paragraph{Reinicio de la búsqueda en cada etapa.} Al comenzar cada etapa de la fase~1, y al
comenzar la fase~2, se restablecen $x$, $\xbest$, la temperatura ($T\gets T_0$) y los
contadores de estancamiento. En cambio, el estado del mecanismo de selección (los pesos de
las ruletas o la tabla $Q$) \emph{se conserva} entre etapas y entre fases; el único aviso que
reciben las subclases es \fn{AlIniciarFase2}, que el método clásico ignora y el Q-learning
aprovecha para volver a explorar (Sección~\ref{sec:ql-fase2}). Obsérvese que dentro de una
etapa de la fase~1 $\xbest$ es la mejor solución \emph{de esa etapa}, normalmente incompleta;
la mejor solución completa se guarda aparte en $\bar{x}$.

\subsection{Ganancias relativas}
\label{sec:ganancia}

El resultado $o$ de cada iteración incluye dos ganancias, respecto de la mejor solución y de
la actual, calculadas \emph{antes} de actualizar $x$ y $\xbest$. La ganancia de una candidata
$x'$ sobre una referencia $y$ es la mejora relativa en el \emph{primer criterio en que
difieren}, en el mismo orden que impone $f$ (\cod{relativeGain}):
\begin{equation}
G(y,x')=
\begin{cases}
\max\Bigl(0,\ \dfrac{|\Uns(y)|-|\Uns(x')|}{|\Uns(y)|}\Bigr) & \text{si } |\Uns(x')|\neq|\Uns(y)|,\\[10pt]
\max\Bigl(0,\ \dfrac{K(y)-K(x')}{K(y)}\Bigr) & \text{si no, y } K(x')\neq K(y),\\[10pt]
\max\Bigl(0,\ \dfrac{D(y)-D(x')}{D(y)}\Bigr) & \text{en otro caso.}
\end{cases}
\label{eq:ganancia}
\end{equation}
Cada mejora se expresa así en la escala de su propio criterio: colocar un cliente del banco
vale $1/|\Uns|$, ahorrar un vehículo vale $1/K$ y una reducción de distancia vale
$\Delta D/D$. Si la ganancia se midiera sobre $f$, las mejoras de distancia quedarían diluidas
por los términos $10^{4}K$ y $10^{6}|\Uns|$, que no cambian. Una candidata que empeora el
primer criterio en que difiere tiene ganancia $0$ aunque mejore los siguientes. Las ganancias
solo las utiliza el Q-learning (Sección~\ref{sec:recompensa}); el método clásico usa
únicamente las banderas de $o$.

\subsection{Criterio de aceptación}
\label{sec:aceptacion}

Se usa el criterio de recocido simulado~\cite{kirkpatrick1983}. Una candidata que no empeora
la solución actual se acepta siempre; una que la empeora en $\Delta=f(x')-f(x)>0$ se acepta
con probabilidad
\begin{equation}
P(\text{aceptar})=\exp(-\Delta/T).
\label{eq:sa}
\end{equation}
Obsérvese que una candidata de igual costo ($\Delta=0$) se acepta pero no se considera mejora,
lo que permite desplazamientos laterales por mesetas de $f$ sin que cuenten como éxito para
el mecanismo de aprendizaje.

Las candidatas incompletas no reciben un tratamiento especial: se evalúan con $f$ como
cualquier otra. En la fase~1 son la norma, y el orden que induce $f$ hace que la búsqueda
prefiera, por este orden, menos clientes pendientes, menos vehículos y menos distancia. Dejar
un cliente más sin asignar supone $\Delta\ge 10^{6}$ menos lo que se ahorre en los otros
términos, y su probabilidad de aceptación es despreciable a las temperaturas de trabajo; en la
práctica, el banco de la solución actual no crece. A igualdad de clientes pendientes, en
cambio, el recocido acepta empeoramientos de distancia con normalidad, y es ese movimiento el
que reordena las rutas hasta que los clientes del banco encuentran hueco.

\subsection{Esquema de temperatura}

\paragraph{Temperatura inicial.} $T_0$ se calibra una sola vez, respecto de la distancia de la
solución inicial, de forma que un empeoramiento de $w\,D(x_0)$ se acepte con probabilidad
$0.5$:
\begin{equation}
T_0=\frac{-\,w\,D(x_0)}{\ln 0.5}=\frac{0.2}{\ln 2}\,D(x_0)\approx 0.2885\,D(x_0).
\label{eq:t0}
\end{equation}
Con esta definición, la probabilidad de aceptar un empeoramiento $\Delta$ a temperatura $T_0$
es $2^{-\Delta/(w D(x_0))}$. Todas las etapas de la fase~1 y la fase~2 arrancan con esta misma
$T_0$. Conviene notar que $T_0$ escala con la distancia de la instancia, mientras que las
penalizaciones de~\eqref{eq:costo} son fijas: la probabilidad de aceptar a $T_0$ una solución
con un vehículo adicional (reutilizando una ruta que había quedado vacía) es
$2^{-5\cdot 10^{4}/D(x_0)}$, que es despreciable cuando $D(x_0)$ es del orden de unos pocos
miles y deja de serlo cuando $D(x_0)$ es del orden de decenas de miles. La de aceptar un
cliente más sin asignar es $2^{-5\cdot 10^{6}/D(x_0)}$, despreciable en todos los tamaños de
instancia considerados.

\paragraph{Enfriamiento.} Tras cada iteración $T\gets c\,T$ con $c=0.9995$. La temperatura
se reduce a la mitad cada $\ln 0.5/\ln c\approx 1386$ iteraciones y, sin recalentamientos,
al cabo de las $25\,000$ iteraciones de la fase~2 valdría
$c^{25000}\,T_0\approx 3.7\cdot 10^{-6}\,T_0$, es decir, un descenso prácticamente puro.

\paragraph{Reinicio por estancamiento.} Si transcurren $L_{\mathrm{est}}=1500$ iteraciones
consecutivas sin un nuevo mejor, la búsqueda vuelve a la mejor solución conocida
($x\gets\xbest$) y la temperatura se fija en $\rho\,T_0=0.5\,T_0$. El reinicio combina dos
efectos: \emph{intensificación}, porque descarta la deriva acumulada y retoma el mejor punto
conocido, y \emph{diversificación}, porque el recalentamiento vuelve a permitir empeoramientos.
Una consecuencia cuantitativa es que, en un periodo prolongado de estancamiento, la
temperatura oscila de forma periódica dentro del intervalo
\[
\bigl[\rho\,c^{L_{\mathrm{est}}}\,T_0,\ \rho\,T_0\bigr]\approx[\,0.236\,T_0,\ 0.5\,T_0\,],
\]
y nunca llega a congelarse. La temperatura solo desciende por debajo de esa banda mientras se
sigan encontrando nuevos mejores con una separación menor que $L_{\mathrm{est}}$ iteraciones.
El mecanismo actúa igual en las dos fases. En la fase~1, «nuevo mejor» es cualquier descenso
de $f$ dentro de la etapa, incluidas las mejoras de distancia a igualdad de banco; por eso el
abandono de un intento no se decide con este contador sino con uno propio
($\textit{sinBanco}$), que solo atiende al número de clientes pendientes.

El contador de estancamiento se actualiza \emph{antes} de \fn{Aprender}, pero el resultado $o$
que recibe el mecanismo de aprendizaje es el de la iteración recién evaluada y no se ve
alterado por el reinicio. Del mismo modo, el cambio de etapa o de fase se decide
\emph{después} de \fn{Aprender}, de manera que la iteración que completa una solución en la
fase~1 se aprende como cualquier otra.

\subsection{Catálogo de operadores y grado de destrucción}
\label{sec:catalogo}

El Algoritmo~\ref{alg:aplicar} muestra cómo se aplica el par elegido. Hay $6$ operadores de
destrucción y $5$ de reparación (Tabla~\ref{tab:catalogo}), es decir, $30$ pares posibles.

\begin{algorithm}[H]
\caption{\fn{AplicarOperadores} (\cod{applyOperators})}
\label{alg:aplicar}
\Entrada{solución $x'$, índice de destrucción $j\in\{0,\dots,5\}$, índice de reparación $k\in\{0,\dots,4\}$}
$q_{\min}\gets\max\bigl(4,\ \min(30,\ \lfloor 0.10\,n\rfloor)\bigr)$\;
$q_{\max}\gets\max\bigl(q_{\min}+1,\ \min(60,\ \lfloor 0.40\,n\rfloor)\bigr)$\;
$q\sim\unifint{q_{\min}}{q_{\max}}$\;
$\textit{des}_j(x', q)$\;
$\textit{ruido}\gets$ verdadero con probabilidad $0.5$\;
$\textit{rep}_k(x', \textit{ruido})$\;
\end{algorithm}

\begin{table}[H]
\centering
\small
\caption{Catálogo de operadores, en el orden de sus índices en el código.}
\label{tab:catalogo}
\begin{tabularx}{\textwidth}{@{}c l L@{}}
\toprule
Índice & Operador & Criterio \\
\midrule
\multicolumn{3}{@{}l}{\emph{Destrucción}}\\
0 & \cod{randomRemoval} & Clientes al azar \\
1 & \cod{routeRemoval} & Rutas completas al azar \\
2 & \cod{worstRemoval} & Clientes cuyo retiro ahorra más distancia (aleatorizado) \\
3 & \cod{shawRemoval} & Clientes parecidos entre sí en espacio, tiempo y demanda \\
4 & \cod{clusterRemoval} & Uno de los dos grupos geográficos de una ruta, y así en rutas vecinas \\
5 & \cod{timeWindowRemoval} & Clientes con las ventanas de tiempo más estrechas \\
\midrule
\multicolumn{3}{@{}l}{\emph{Reparación} (cada una, con o sin ruido en los costos)}\\
0 & \cod{greedyInsertion} & Inserción de menor costo \\
1 & \cod{regret2Insertion} & Mayor arrepentimiento entre las dos mejores posiciones \\
2 & \cod{regret3Insertion} & Mayor arrepentimiento entre las tres mejores rutas \\
3 & \cod{regret4Insertion} & Mayor arrepentimiento entre las cuatro mejores rutas \\
4 & \cod{regretMInsertion} & Mayor arrepentimiento sobre todas las rutas \\
\bottomrule
\end{tabularx}
\end{table}

\paragraph{Grado de destrucción.} Se sortea de forma uniforme entre el 10\,\% y el 40\,\% de
los clientes, con los topes absolutos de Pisinger y Ropke~\cite{pisinger2007}:
\[
q\in\bigl[\min(0.1\,n,\ 30),\ \min(0.4\,n,\ 60)\bigr],
\]
es decir, $q\in[10,40]$ para $n=100$, $q\in[20,60]$ para $n=200$ y $q\in[30,60]$ para todo
$n\ge 300$. En las instancias de $400$, $600$ y $800$ clientes el vecindario tiene, por tanto,
el mismo tamaño absoluto y representa una fracción decreciente de la solución (entre el
3.75\,\% y el 7.5\,\% para $n=800$). El tope mantiene acotado el costo de una iteración, que
está dominado por la reparación (Sección~\ref{sec:reparacion}), y es lo que hace asumible un
presupuesto de hasta $2\,T_{\max}=50\,000$ iteraciones en las instancias grandes. En la fase~1
la reparación debe reinsertar, además de los $q$ clientes extraídos, los que ya estaban en el
banco.

\paragraph{Ruido.} En cada iteración se decide con probabilidad $0.5$, de forma independiente
del par elegido, si la reparación trabaja con los costos de inserción exactos o con costos
perturbados (Sección~\ref{sec:ruido}). El ruido no es un operador aparte ni una decisión del
mecanismo de aprendizaje: cada operador de reparación se aplica, en promedio, la mitad de las
veces en cada modalidad, y el crédito que recibe mezcla ambas.
